Calibration free, with each ear referenced to itself as two earbuds must be, dry contacts in the ear canal and concha carried no attention that a decoder trained on other people could read: \earInNn{} of one minute decisions with the learned decoder and \earInLin{} with the linear one, at a coin flip. Wet contacts around the same ears, at the same moments, gave \earAroundNn{} (\earAroundNnFive{} at five minutes), and putting the two sets together gave \earBothNn{}, less than around the ear alone: for people the decoder has never seen, the in-ear channels are noise to fit. The data set's own analysis, trained on each person, put dry in-ear decoding at 61.1\% and found that dropping the calibration cost it 11.7 points~\citep{geirnaert2025}, which is where this calibration free result sits. Two confounds remain, and the data cannot separate them: the in-ear contacts were dry and the around-ear ones wet, and the in-ear ones were fewer and closer together. Neither changes what an earbud that has to work out of the box needs, which is contacts that reach outside the canal. The specification that accompanies this work now requires at least two of the four sensing contacts per ear outside it.
